\section{Introducción}
\label{sec:intro}

% P1 — Definición + fuente
La batimetría es la medición y representación de la profundidad y el relieve
del fondo de un cuerpo de agua respecto de un nivel de referencia vertical, y
constituye en esencia el equivalente sumergido de la topografía terrestre
\cite{wolfl2019, talley2011}. Describe tanto la profundidad local como la
forma del lecho de océanos, estuarios, ríos y lagos.

% P2 — Por qué importa
La batimetría es un insumo fundamental de la hidráulica, la hidrología y las
ciencias costeras y oceánicas. El lecho impone la condición de borde inferior
sobre la que se desarrolla todo el flujo, de modo que la propagación de
inundaciones y tsunamis, la circulación de mareas y corrientes, el transporte
de sedimentos y la seguridad de la navegación dependen de ella. En la práctica,
la fidelidad de cualquier modelo hidrodinámico está acotada por la fidelidad de
la batimetría que se le suministra.

% P3 — Cómo funciona el sonar multihaz
El levantamiento batimétrico directo se realiza principalmente mediante dos
tecnologías. La primera es la ecosonda multihaz, instalada en una embarcación,
que emite un abanico de haces acústicos transversales a la trayectoria y mide
el tiempo de viaje de ida y vuelta de cada eco reflejado por el fondo. Conocida
la velocidad del sonido en la columna de agua, ese tiempo y el ángulo del haz
determinan la profundidad y la posición de cada sondaje sobre una franja ancha,
que se cubre de forma continua a medida que la embarcación avanza
\cite{lurton2010}.

% P4 — Cómo funciona el LiDAR batimétrico
La segunda es el LiDAR batimétrico aerotransportado, que emite pulsos de láser
de longitud de onda verde capaces de penetrar la columna de agua. La
profundidad se obtiene de la diferencia temporal entre el retorno reflejado por
la superficie libre y el reflejado por el fondo, lo que permite cartografiar con
rapidez aguas someras y claras desde una aeronave \cite{guenther2000,
szafarczyk2023}.

% P5 — Brecha de cobertura / costo (citada, redacción conservadora)
Ambas técnicas son precisas pero lentas, costosas y logísticamente limitadas:
requieren plataformas dedicadas, condiciones operativas favorables y campañas
que deben repetirse, porque el lecho evoluciona por sedimentación y erosión.
Como resultado, solo una fracción del fondo oceánico mundial está cartografiada
con métodos acústicos modernos, del orden de una cuarta parte a mediados de la
década de 2020 frente a apenas unos pocos puntos porcentuales una década antes,
y la cobertura por sonar multihaz de alta resolución es considerablemente menor
\cite{wolfl2019}. Esta brecha persistente motiva el desarrollo de métodos
indirectos que infieran el fondo a partir de observaciones más accesibles.

% P6 — Genérico, sin nombrar Chile
%% IEEE-CUT: en la versión IEEE este párrafo se reduce a media frase ("regiones costeras con datos escasos").
El déficit es más severo donde los levantamientos son escasos, gruesos o están
desactualizados: incluso un sitio ya sondeado pierde vigencia porque el lecho
evoluciona por sedimentación y erosión, de modo que la necesidad no es solo
cartografiar por primera vez, sino actualizar batimetría obsoleta o de baja
resolución, a menudo en tramos de difícil acceso. Frente a ese costo, el
monitoreo del nivel del agua es barato, continuo y permanente, y no requiere un
levantamiento batimétrico para operar: los registros de nivel mareal y de nivel
y caudal fluvial se acumulan de forma sostenida justo allí donde falta una
batimetría distribuida de alta resolución. Las mediciones de velocidad
superficial aportan una restricción más fuerte, pero son más escasas y suelen
asociarse a despliegues dedicados; aun así pueden obtenerse de forma
independiente, por ejemplo mediante boyas a la deriva, velocimetría por video o
radar de alta frecuencia, por lo que el valor central del enfoque descansa sobre
la señal de nivel y su variación temporal. Surge así la oportunidad de un método
inverso que actualice batimetría obsoleta o gruesa, incluso en tramos
inaccesibles, a partir de datos de superficie que ya se registran de forma
continua, en lugar de encargar nuevas campañas.

% P7 — Idea del problema inverso (lenguaje natural, sin notación)
La idea central de este trabajo es, en lugar de medir el fondo directamente,
inferirlo de la huella que deja sobre el agua. La forma del lecho controla cómo
se elevan y deforman la superficie libre y cómo se aceleran las corrientes; por
lo tanto, las observaciones de la superficie y del flujo contienen información
sobre el fondo no observado. Recuperar el lecho a partir de esas observaciones
es un problema inverso, mal planteado en general, que se vuelve tratable cuando
se lo restringe con la física que gobierna el agua somera.

% P8 — El enfoque (conceptual; notación diferida)
Para resolver ese problema inverso, esta tesis adopta las \textbf{redes
neuronales informadas por física} (PINNs) \cite{raissi2019}: redes entrenadas
para ser consistentes, de forma simultánea, con las observaciones de superficie
y con las ecuaciones de aguas someras que gobiernan el flujo. Bajo esa doble
restricción, el fondo desconocido emerge como el campo que explica los datos y
satisface las ecuaciones a la vez. El enfoque emplea, conceptualmente, un
modelo para el campo de flujo y otro independiente para el lecho. La
formulación matemática de las ecuaciones de aguas someras y toda la notación se
introducen en la Sección~\ref{sec:problema}; la arquitectura y la
función de pérdida, en la Sección~\ref{sec:method}.
